\section{An extension and reality gate for ordinary orthogonal levels}
\label{sec:orthogonal-lowlevel-reality}

\paragraph{Status and scope.}
This is a new author derivation.  The complete proof, invoked source
statements and exact author certificate have passed the author check
and complete independent mathematical reviews by /root and
/root/g34\_eightfolds. The peer also supplied a separately authored
exact supplement; all actual scopes and timings are recorded.  The new orbit argument is separate
from the published character estimates in D-OL-01.  It gives a strict
reduction of the remaining unbounded natural dimension case, without
asserting that any surviving case is a source.

Write $O(V)$ for the full isometry group of a nondegenerate quadratic
space over $\mathbf F_q$, and put $S=\Omega(V)=[O(V),O(V)]$.
For odd $q$, the natural dimension $d$ may be even or odd; for even $q$
we use a nondefective even-dimensional quadratic space.  The
orthogonal level $\ell$ is the level of Definition~3.2(ii) of
\cite{OL:GLT2024}.  In particular, Lemma~3.3 there identifies its
constituents on $S$ with those of powers of the permutation character
$\tau_V$ on the vector set $V$.

\begin{lemma}[D-OL-04: extension from a short tuple action]
\label{D-OL-04}
Assume $d\ge10$.  If an irreducible complex character $\rho$ of $S$
has level $\ell$ and
\[
                         d\ge4\ell+2,                         \tag{OL4.1}
\]
then $\rho$ extends to an irreducible character of $O(V)$ and is
self-dual.  Consequently every nonselfdual irreducible character of
$S$ satisfies
\[
                         d\le4\ell+1.                         \tag{OL4.2}
\]
The extension assertion concerns an actual representation, not just
invariance of its character under outer automorphisms.
\end{lemma}

\begin{proof}
Let $\mathcal M=\mathbf C[V^\ell]$.  We prove that $O(V)$ and $S$
have identical orbits on $V^{2\ell}$, which is the set indexing the
matrix entries of endomorphisms of $\mathcal M$.

Fix a tuple of $2\ell$ vectors, and let $R$ be its span, of dimension
$t\le2\ell$.  For the nondegenerate polar form on $V$, the radical
of $R^\perp$ is $R\cap R^\perp$.  Its dimension is at most $t$.
Thus the induced nondegenerate polar space after taking out the radical
has dimension at least $d-2t\ge d-4\ell\ge2$.  Choose a vector-space
complement $U$ to this radical inside $R^\perp$.  The polar form
restricted to $U$ is nondegenerate; in characteristic two this means
that the restricted quadratic space is nondefective.  Hence
\[
 V=U\perp U^\perp,\qquad R\subseteq U^\perp,\qquad \dim U\ge2.
                                                               \tag{OL4.3}
\]
Every element of $O(U)$ extended by the identity on $U^\perp$ fixes
the chosen tuple pointwise.

We check that this embedded $O(U)$ surjects onto $O(V)/S$.
For odd $q$, determinant and spinor norm identify this quotient with
$\{\pm1\}\times\mathbf F_q^*/(\mathbf F_q^*)^2$;
see the precise subgroup descriptions in
\cite[Section~2, pp.~41--43]{OL:KTV2023}.
A nondegenerate binary subspace of $U$ has a diagonal quadratic form
$ax^2+by^2$ with $ab\ne0$.  It represents every nonzero $c\in\mathbf F_q$:
the two subsets $a\{x^2:x\in\mathbf F_q\}$ and
$c-b\{y^2:y\in\mathbf F_q\}$ each have $(q+1)/2$ elements, so they
intersect.  Reflections in vectors of square and nonsquare nonzero
norm therefore belong to $O(U)$ and have quotient classes
$(-1,1)$ and $(-1,\eta)$, respectively, where $\eta$ is the
nontrivial square class.  These classes generate the entire quotient.
The same conclusion holds for either type of $U$ and for odd $d$.

For even $q$, the quotient $O(V)/S$ has order two
\cite[Section~2, p.~42]{OL:KTV2023}.  There is a vector $v\in U$
with $Q(v)\ne0$, since otherwise the polar form on $U$ would vanish.
The symmetry
\[
              s_v(x)=x+\frac{B(x,v)}{Q(v)}v                    \tag{OL4.4}
\]
is an isometry supported on $U$ and has $\operatorname{rank}(s_v-1)=1$.
It is outside $S$: the even-dimensional fixed-space property on $S$
used in the proof of \cite[Lemma~3.3, p.~160]{OL:GLT2024} would
otherwise give even rank.  Thus $O(U)$ surjects onto the quotient
also in characteristic two.

For any $g\in O(V)$, choose a tuple stabilizer $h\in O(U)$ with the
same class as $g$ in $O(V)/S$.  Then $gh^{-1}\in S$ sends the tuple
to exactly its image under $g$.  This proves equality of the two orbit
partitions.  Since an equivariant matrix on a permutation module is
exactly a matrix constant on the corresponding pair orbits, we obtain
an equality of algebras, not only an equality of their dimensions:
\[
       \operatorname{End}_{S}(\mathcal M)
       =\operatorname{End}_{O(V)}(\mathcal M).                 \tag{OL4.5}
\]

Decompose $\mathcal M$ into irreducible $O(V)$-summands over
$\mathbf C$.  For an irreducible summand $W$, every
$S$-endomorphism of $W$ extends by zero on an $O(V)$-stable complement
to an $S$-endomorphism of $\mathcal M$.  Equation (OL4.5) makes it
$O(V)$-equivariant, so its restriction to $W$ is scalar by Schur's
lemma.  Complete reducibility then shows that $W|_S$ is irreducible.
The same argument applied to homomorphisms between two summands
shows that distinct $O(V)$-types restrict to distinct $S$-types.

On $S$ one has $\tau_V+\zeta_V=2\tau_V$ by
\cite[Lemma~3.3, p.~160]{OL:GLT2024}.  Thus a character of level
$\ell$ occurs in $\mathcal M|_S$, and the preceding paragraph
supplies its irreducible extension to $O(V)$.
Finally the full finite orthogonal group $O(V)$ is real in both
characteristics: every element is a product of two involutions.
For odd characteristic this is Wonenburger's theorem, restated
precisely in \cite[Theorem~1, p.~44]{OL:KTV2023}; for characteristic
two the same published paper states the Ellers--Nolte and Gow theorem
explicitly on p.~40.  These known results imply that every irreducible
character of $O(V)$ is real-valued.  Its representation is therefore
isomorphic to its dual.  Restricting the extension proves the
self-duality of $\rho$.  No assertion about its Frobenius--Schur sign
is needed.  The contrapositive gives (OL4.2).
\end{proof}

\begin{theorem}[D-OL-05: the remaining unbounded orthogonal case]
\label{D-OL-05}
Retain the actual six-dimensional multiplicity variation and
rank-nine period-map setting of D-OL-02--03.  Suppose the inducing
image $J$ has a normal irreducibly acting central product of
quasisimple orthogonal factors, permuted transitively by $J$.
Let $d$ and $q$ be their common natural dimension and field size,
and let the original irreducible representation have imaginary
quadratic character field.  Put $J_0=10^{30}$ as in (OL.1).

Every nonselfdual ordinary factor character satisfies
\[
                         d\le4J_0-3.                         \tag{OL4.6}
\]
More generally, if
\[
                         d\ge4J_0-2,                         \tag{OL4.7}
\]
then all factor characters are ordinary, have level less than
$J_0$, extend irreducibly to the full orthogonal group, and are
self-dual.  The inducing image $J$ is essentially self-dual with
nontrivial multiplier.  In addition,
\[
                 q\in\{2,3,4,5,7,8,9,11,13\}.                \tag{OL4.8}
\]
In particular every transitive orthogonal-core source with
$q\ge16$ has $d\le4J_0-3$.  All conclusions allow arbitrary
compatible scalar, diagonal, graph and field normalizers, and
arbitrary transitive factor permutations.  They do not require
equality between the inducing character field and the original
imaginary quadratic field.
\end{theorem}

\begin{proof}
D-OL-03 gives $\ell<J_0$ for every ordinary factor character.
Equation (OL4.2) immediately gives (OL4.6).  Under (OL4.7), the
spin and nonextension cutoffs of D-OL-03 already force the factor
characters to be ordinary.  Since $\ell\le J_0-1$, one has
$4\ell+2\le4J_0-2\le d$.  D-OL-04 therefore gives all the stated
extension and self-duality assertions.

Let $N\triangleleft J$ be the indicated irreducible central
product and let $e$ be the degree of the inducing representation
$\tau$.  The irreducible $N$-module is a tensor product of the
irreducible factor modules.  Tensoring their nondegenerate invariant
bilinear forms yields such a form for $N$.  Its invariant-form line
is one-dimensional by Schur's lemma, and normality makes $J$
preserve that line.  Thus $J$ has a multiplier character
$\mu:J\to\mathbf C^*$, whether the resulting form is symmetric or
alternating.  This also handles elements permuting the factors;
no separate extension or compatibility assumption is needed.

If $\mu$ were trivial, $\tau$ would be self-dual and its character
would be real-valued.  But the original character is induced from
$\tau$, so the induction formula puts the original character field
inside the character field of $\tau$ at the given complex embedding.
The former is imaginary quadratic and cannot lie in $\mathbf R$.
Consequently $\mu$ is nontrivial.

Suppose now that $q\ge16$.  The published uniform character ratio
bound used in D-BI-18 is
\[
 c(q)=\max\{1/(\sqrt q-1),9/q\}\le9/16.                       \tag{OL4.9}
\]
Indeed the first term is at most $1/3$ and the second at most
$9/16$.  This is \cite[Theorem~4.1(2), p.~141]{BI:DHP}, with its
quasisimple Lie-type and nonalternating quotient hypotheses
satisfied here.  Tensoring the factor characters shows that every
nonscalar element of $N$ has character ratio at most $9/16$.
The complete commutator argument of D-BI-18, including the case
of a nonidentity scalar commutator, therefore gives
\[
             \nu(g)\ge\frac{7e}{64}\quad
                     \text{for every nonscalar }g\in J.      \tag{OL4.10}
\]
This conclusion uses only normality and irreducibility of $N$;
it includes all outer normalizers and all factor permutations.
A genuine nonidentity scalar inertia has actual moving rank $e$,
so it also satisfies the needed lower bound on moving rank.

The actual intermediate Hurwitz factorization gives genus zero,
at least twelve genuine branch points, and total moving rank
$2e+6$ as in D-BI-18.  The inertias generate $J$, have product one,
and $\mu$ is nontrivial, so at least two have multiplier different
from one.  For such an inertia the $1$-eigenspace pairs only with
the $\mu(g)$-eigenspace under the nondegenerate form, and hence
has dimension at most $e/2$.  The two corresponding moving ranks
are at least $e/2$.  The other at least ten inertias satisfy (OL4.10).
Consequently
\[
  2e+6\ge e+10\frac{7e}{64}=\frac{67e}{32},\qquad e\le64.
                                                               \tag{OL4.11}
\]
The factor characters are nontrivial.  The $k=1$ case of
\cite[Theorem~3.7, pp.~165--167]{OL:GLT2024} gives
\[
                         a\ge q^{d-4}(q-1),                  \tag{OL4.12}
\]
where $a$ is the degree of any one factor: below that threshold
the character has level zero and is trivial.  Here $d\ge10$ and
$q\ge16$, so $a\ge16^6\cdot15>64$.  Since $e$ is the product of
the factor degrees, $e\ge a$, contradicting (OL4.11).
Thus $q<16$, and listing the prime powers below sixteen gives
exactly (OL4.8).

All uses of the branch count are necessary consequences of the
actual period-map factorization.  We have not inferred dominance
from a dimension inequality, nor ruled out the ordinary self-dual
bounded-level cases in (OL4.8).
\end{proof}
